\subsection{Ordinary nullification on the countable H-space gate}
\label{sec:A4-hspace}

The countable reduction of A4.2 can be restricted further while
retaining countability.  The point beyond the loop--suspension
replacement is that ordinary and reversible $BC_p$-acyclicity agree
on the resulting H-spaces.  The following criterion leaves the
existence question open.

\begin{claim}{A4.5}{Proved}{A proper countable H-space domain with an unsuspended nullification test}
\label{A45}
Let $p$ be a prime and write $P=P_{BC_p}$.  If $H$ is a simply
connected H-space with $p$-primary reduced integral homology, then
\[
 P(\Sigma H)\simeq *
 \quad\Longleftrightarrow\quad PH\simeq *
 \quad\Longleftrightarrow\quad
 H\in\mathcal A_{\min\operatorname{Rev}(BC_p)}.
 \tag{A4.5.1}
\]
At $p=2$, put $W_0=\min\operatorname{Rev}(BC_2\to*)$ and
$W_*=W_S\cap W_{\mathbb Z[1/2]}$.  Then $W_0=W_*$ if and only if
\[
 P_{BC_2}H\simeq *
 \tag{A4.5.2}
\]
for every \emph{countable simply connected H-space} $H$ with
$2$-primary reduced integral homology and
$\Map_*(H,S^n_{(2)})\simeq *$ for every $n\geq2$.
It even suffices to test the spaces
\[
 H=J(\Sigma X)\simeq\Omega\Sigma^2X,
 \tag{A4.5.3}
\]
where $X$ runs over the countable connected spaces of A4.2.
The H-space domain is a proper subclass of that previous domain,
even after restricting the latter to simply connected spaces.

Equivalently, $W_0\subsetneq W_*$ holds if and only if there exists
a noncontractible simply connected $2$-primary H-space $Z$ which
is sphere-null as a source and $BC_2$-null as a target.  No
countability assertion is made about $Z$ in this last formulation.
\end{claim}
\begin{proof}
We first prove \textup{(A4.5.1)}.  The equivalence with reversible
asphericity is A34.2, once the first equivalence is proved.  Ordinary
$BC_p$-acyclic spaces are closed under suspension, so only
$P\Sigma H\simeq *\Rightarrow PH\simeq *$ needs proof.
Put $Z=PH$.

\paragraph{Preservation of the required properties.}
The construction following \citet[Theorem~1.1, p.~34]{ABousfield1996}
attaches mapping cones on maps from $BC_p$ and its nonnegative
suspensions, and takes a suitable transfinite homotopy colimit.
Each attaching domain is connected and has zero reduced homology
over $\mathbb Z[1/p]$.  Coning such a map into a simply connected
space preserves simple connectivity by van Kampen.  The cofiber
homology sequence shows that it preserves homology over
$\mathbb Z[1/p]$.  These conclusions persist through the indicated
homotopy colimits and CW replacements.  Thus $Z$ remains simply
connected and primary.  No bound on the torsion exponents is used.
Nullification preserves finite products, as stated explicitly in
the proof of \citet[Theorem~1.7, p.~35]{ABousfield1996}.
Applying $P$ to the multiplication of $H$ and its unit homotopies,
and using $P(H\times H)\simeq PH\times PH$, makes $Z$ an H-space.
No preservation of co-H structures is asserted.

\paragraph{The symmetric-product retraction.}
The space $Z$ is $BC_p$-null, and so is $\Omega Z$, because pointed
mapping spaces commute with the homotopy limit defining loops.
The universal property of $H\to PH$ and the assumed acyclicity
of $\Sigma H$ give equivalences of whole pointed mapping spaces
\[
 \Map_*(\Sigma Z,Z)
 \simeq\Map_*(Z,\Omega Z)
 \simeq\Map_*(H,\Omega Z)
 \simeq\Map_*(\Sigma H,Z)
 \simeq *.
 \tag{A4.5.4}
\]
In particular this controls components as well as positive
homotopy groups.  It does not identify $P\Sigma H$ with $\Sigma PH$.
The exact input is \citet[Key Lemma~1.6, p.~35]{ABousfield1996}:
if $V$ is connected and $T$ is a connected $\Sigma V$-null H-space,
the inclusion $V\to SP^\infty V$ induces an equivalence
\[
 \Map_*(SP^\infty V,T)\xrightarrow{\simeq}\Map_*(V,T).
 \tag{A4.5.5}
\]
Apply it with $V=T=Z$, using \textup{(A4.5.4)}.  The identity of
$Z$ extends along $Z\to SP^\infty Z$, so $Z$ is a homotopy
retract of $SP^\infty Z$.  The Dold--Thom description on the same
source page makes the latter a simply connected GEM with
$p$-primary homotopy groups $H_n(Z;\mathbb Z)$.  By A34.2,
using \citet[Lemma~8.13]{ABousfield1997}, every such GEM is
ordinarily $BC_p$-acyclic.  Acyclic classes are closed under
retracts, so $Z$ is both $BC_p$-null and $BC_p$-acyclic.
Consequently $Z\simeq *$, proving \textup{(A4.5.1)}.

\paragraph{The countable replacement.}
If $W_0=W_*$, every H-space in the stated domain is
$W_0$-aspherical; \textup{(A4.5.1)} proves
\textup{(A4.5.2)}.  Conversely let $X$ be any countable
connected source in A4.2, set $A=\Sigma X$, and form $H=J(A)$.
The reduced-product construction of
\citet[Section~1 and Theorem~5.6]{A4James1955} provides a
countable CW model and a natural equivalence
$J(A)\simeq\Omega\Sigma A$.  The proof of Theorem~5.6 on
pp.~184--187 includes the passage from finite subcomplexes
to countable CW complexes.  Its hypotheses hold because
$A$ is countable and connected.

We check all properties needed here.  The reduced-product cells
are represented by finite words of cells of $A$, so there are
countably many of them; concatenation supplies the multiplication.
The space $\Sigma^2X$ is $2$-connected, so $H\simeq\Omega\Sigma^2X$
is simply connected.  The James filtration has
$J_0(A)=*$ and cofiber sequences
\[
       J_{r-1}(A)\longrightarrow J_r(A)\longrightarrow A^{\wedge r}
       \qquad(r\geq1).
 \tag{A4.5.6}
\]
The $W_*$-aspherical class is closed under suspension, smash with
pointed spaces, cofiber extensions, and pointed homotopy colimits.
Starting from $X$ it therefore contains $A$, every
$A^{\wedge r}$, each $J_r(A)$, and their sequential homotopy
colimit $H$.  In particular $H$ is primary and sphere-null.
This uses unstable cofiber sequences, without a stable
inverse-limit comparison or a suspension splitting.

Finally the suspension--loop adjunction triangle gives
\[
 \Sigma A\xrightarrow{\Sigma\eta_A}\Sigma\Omega\Sigma A
       \xrightarrow{\epsilon_{\Sigma A}}\Sigma A,
 \qquad
 \epsilon_{\Sigma A}\circ\Sigma\eta_A\simeq\id_{\Sigma A}.
 \tag{A4.5.7}
\]
Thus $\Sigma^2X$ is a homotopy retract of $\Sigma H$.
If \textup{(A4.5.2)} holds for this $H$, ordinary acyclicity
puts $H$ in $\mathcal A_{W_0}$.  Suspension, retract closure,
and reversible desuspension twice then put $X$ in
$\mathcal A_{W_0}$.  A4.2 implies $W_*=W_0$.
Contrapositively, every countable separator $X$ produces a
countable H-space separator $H$ with $PH\not\simeq *$.
No countability of $PH$ is needed.

\paragraph{Strictness of the domain restriction.}
The space $V=\Sigma BC_p$ is countable, simply connected,
primary and sphere-null, but cannot admit an H-space structure.
Indeed choose nonzero classes
\[
 a\in H^2(V;\mathbb F_p),\qquad b\in H^3(V;\mathbb F_p).
\]
They exist by the suspension isomorphism and the nonzero
degree-one and degree-two cohomology of $BC_p$.
Since $H^1(V;\mathbb F_p)=0$, the pullback of a unital
multiplication would make both $a$ and $b$ primitive: the only
possible positive-degree decompositions of degrees two and three
involve degree one.  Positive-degree cup products on a suspension
vanish, so $ab=0$.  Multiplicativity of the pullback would give
\[
 0=\mu^*(ab)=a\otimes b+(-1)^{2\cdot3}b\otimes a.
 \tag{A4.5.8}
\]
The summands are nonzero and lie in distinct K\"unneth bidegrees
$(2,3)$ and $(3,2)$, a contradiction, also at $p=2$.
The groups are finite dimensional in each degree, so this
K\"unneth calculation has no completion issue.
This verifies a proper restriction of the quantified domain.
The example $\Sigma BC_p$ itself is in $\mathcal A_{W_0}$;
it is not a separator.

\paragraph{The null H-space formulation.}
For a countable H-space separator $H$, take $Z=PH$.
The first part of the proof makes $Z$ simply connected,
primary and an H-space, and \textup{(A4.5.1)} makes it
noncontractible.  It is $BC_p$-null by definition.
Each sphere target $S$ in the sphere-null test is $BC_p$-null
by the Sullivan theorem and its localized/completed comparison
already used in A4.2.  Thus $\Omega^jS$ is $BC_p$-null for
every $j\geq0$, and
\[
 \Map_*(\Sigma^jZ,S)
 \simeq\Map_*(Z,\Omega^jS)
 \simeq\Map_*(H,\Omega^jS)
 \simeq\Map_*(\Sigma^jH,S)
 \simeq *.
 \tag{A4.5.9}
\]
This proves sphere-nullity with all positive suspensions and
components checked.  Conversely any noncontractible $Z$ with
the stated properties satisfies $PZ\simeq Z\not\simeq *$.
Equation \textup{(A4.5.1)} excludes it from $W_0$, while its
primary sphere-nullity puts it in $W_*$.  This proves the last
equivalence without asserting that nullification preserves a
countable number of cells.
\end{proof}

\paragraph{Scope of the reduction.}
Replacing $X$ by $\Omega\Sigma^2X$ alone is a normal-form
observation for reversible classes.  The further input in A4.5
is the ordinary-versus-reversible comparison
\textup{(A4.5.1)}: on the retained countable H-space domain,
the one initial suspension in A34.2 can be removed from the
ordinary-nullification test.  The domain is proper by
\textup{(A4.5.8)}.  This criterion decides neither whether all
its H-spaces are ordinarily $BC_2$-acyclic, nor the existence
of a third interval element or the size of the full stable-part
fibre.  No priority claim is made.

\begin{claim}{A4.6}{Proved}{Finite Steenrod-algebra generation cannot supply a separator}
\label{A46}
Let $p$ be any prime and let $H$ be a simply connected H-space
with $p$-primary reduced integral homology.  Suppose that
\[
 \Map_*(H,S^n_{(p)})\simeq *\quad(n\geq2)
\]
and that $H^*(H;\mathbb F_p)$ is finitely generated as an algebra
over the Steenrod algebra $\mathcal A_p$.  Then
\[
                         P_{BC_p}H\simeq *.
 \tag{A4.6.1}
\]
Consequently the countable H-space test in A4.5 need only be
checked on H-spaces whose mod-$2$ cohomology is \emph{not}
finitely generated as an unstable algebra.  This result does not
replace finite generation by degreewise finite-dimensional
cohomology.
\end{claim}
\begin{proof}
Write $Z=P_{BC_p}H$.  The precise primary input is the published
\citet[Theorem~7.2]{A4CCSDeconstruct2007}: for a connected H-space
whose mod-$p$ cohomology is finitely generated as an algebra over
$\mathcal A_p$, its $BC_p$-nullification has finite total mod-$p$
cohomology.  This theorem does not impose $p$-completeness on $H$.
Thus there is an integer $N\geq1$ with
\[
                         H_n(Z;\mathbb F_p)=0\qquad(n>N).
 \tag{A4.6.2}
\]
By the preservation arguments in A4.5, $Z$ is simply connected,
primary and sphere-null.  The homology universal coefficient
sequence gives surjections
\[
 H_{k+1}(Z;\mathbb F_p)\longrightarrow
 \operatorname{Tor}^{\mathbb Z}_1
       (H_k(Z;\mathbb Z),\mathbb F_p)
       =H_k(Z;\mathbb Z)[p].
 \tag{A4.6.3}
\]
For $k\geq N$ the left side is zero.  A nonzero $p$-primary
torsion group has an element of order $p$, so
$H_k(Z;\mathbb Z)=0$ in all those degrees.  Hence the reduced
integral homology of $Z$ is bounded above.  This argument does
not require the integral groups to be finitely generated and
includes divisible $p$-primary groups.

Suppose that this reduced homology is nonzero.  Let $d$ be its
greatest nonzero degree and put $A=H_d(Z;\mathbb Z)$.
Simple connectivity gives $d\geq2$.  The cohomology universal
coefficient theorem, and the short exact sequence
\[
 0\longrightarrow\mathbb Z_{(p)}\longrightarrow\mathbb Q
 \longrightarrow\mathbb Q/\mathbb Z_{(p)}\longrightarrow0,
\]
give
\[
 H^{d+1}(Z;\mathbb Z_{(p)})
 \cong\operatorname{Ext}^1_{\mathbb Z}(A,\mathbb Z_{(p)})
 \cong\operatorname{Hom}(A,\mathbb Q/\mathbb Z_{(p)})\ne0.
 \tag{A4.6.4}
\]
Indeed $\operatorname{Hom}(A,\mathbb Q)=0$ because $A$ is
torsion, and $\operatorname{Ext}^1(A,\mathbb Q)=0$ because
$\mathbb Q$ is injective.  The $p$-primary Pr\"ufer group
$\mathbb Q/\mathbb Z_{(p)}$ admits a nonzero map from a
nonzero cyclic subgroup of $A$, and that map extends to $A$
by injectivity.  This proves the last nonvanishing for arbitrary
$p$-primary $A$.

Choose a nonzero class $c$ in \textup{(A4.6.4)}.  It represents
a pointed map
\[
                 Z\longrightarrow K(\mathbb Z_{(p)},d+1),
\]
the bottom Postnikov section of $S^{d+1}_{(p)}$.
The same universal coefficient theorem gives
\[
 H^r(Z;B)=0\quad\text{for every abelian group }B
                   \text{ and every }r>d+1.
 \tag{A4.6.5}
\]
In the simply connected sphere's Postnikov tower the obstruction
to lifting from the $k$th to the $(k+1)$st stage lies in
$H^{k+2}(Z;\pi_{k+1}S^{d+1}_{(p)})$, for $k\geq d+1$.
All these groups vanish by \textup{(A4.6.5)}.  Choosing lifts
successively in fibration models also supplies their compatibility.
Postnikov convergence therefore gives a map
\[
                         Z\longrightarrow S^{d+1}_{(p)}
 \tag{A4.6.6}
\]
whose bottom Postnikov class is $c$.  It is nonnull, contrary
to sphere-nullity.  Thus $Z$ is integrally acyclic; since it is
simply connected, the homological Whitehead theorem makes it
contractible.  This proves \textup{(A4.6.1)}.
\end{proof}

\paragraph{The remaining finiteness gap.}
The theorem just used is the published Theorem~7.2, whose proof
uses the Andr\'e--Quillen argument in Theorem~6.1 of
\citet{A4CCSDeconstruct2007}; the earlier author manuscript has
different numbering and a superseded supporting proof.
Finite generation as an unstable algebra is stronger than
finite dimensionality in each degree.  Expressing a cohomology
algebra as a union of finitely generated unstable subalgebras
does not realize that union by H-spaces which remain primary
and sphere-null.  No such realization or reduction of all
countable candidates has been established here.  The unrestricted
criterion of A4.5 therefore remains open.

\begin{claim}{A4.7}{Proved}{Finite cohomology and finite Steenrod generation reach only the endpoints}
\label{A47}
Fix a prime $p$, put
$W_0=\min\operatorname{Rev}(BC_p\to*)$ and
$U_p=W_{\mathbb Z[1/p]}$, and call a connected space \emph{primary}
when its reduced integral homology consists of $p$-primary torsion groups.
The following assertions hold.

\begin{enumerate}
\item If $X$ is primary and
$0<\dim_{\mathbb F_p}\widetilde H_*(X;\mathbb F_p)<\infty$,
then the reversible aspherical class generated by $X$ contains a
noncontractible finite simply connected primary complex.
\item If $V$ is a finite connected primary complex with nonzero
reduced integral homology, then
\[
          \min\operatorname{Rev}(BC_p\vee V\to*)=U_p.
          \tag{A4.7.1}
\]
In particular the conclusion remains true with $V$ replaced by
any $X$ satisfying the first assertion.
\item Let $H$ be a simply connected primary H-space whose mod-$p$
cohomology is finitely generated as an algebra over the Steenrod
algebra.  Then
\[
 \min\operatorname{Rev}(BC_p\vee H\to*)=
 \begin{cases}
  W_0,&P_{BC_p}H\simeq *,\\
  U_p,&P_{BC_p}H\not\simeq *.
 \end{cases}
 \tag{A4.7.2}
\]
Consequently any reversible localizer in $[W_0,U_p]$ generated
over $W_0$ by a set or class of such H-spaces is an endpoint.
\end{enumerate}
The second assertion rechecks the supplied REVLOC~D.3.3; its
periodicity inputs and their connectivity bounds are made explicit below.
\end{claim}
\begin{proof}
\par\medskip\noindent\textbf{Extracting a finite primary complex.}\quad
Put $Y=\Sigma X$ and let $p_Y:Y\to Y$ be the $p$-fold sum
of the identity under the suspension co-H structure.  Let
$C=\cofib(p_Y)$.  The space $C$ is simply connected and belongs
to every reversible aspherical class containing $X$.  The integral
homology exact sequence gives
\[
 0\longrightarrow \widetilde H_k(Y;\mathbb Z)/p
 \longrightarrow \widetilde H_k(C;\mathbb Z)
 \longrightarrow \widetilde H_{k-1}(Y;\mathbb Z)[p]
 \longrightarrow0.
 \tag{A4.7.3}
\]
By the universal coefficient sequence over $\mathbb F_p$, both
outer groups are finite dimensional $\mathbb F_p$-vector spaces,
and they vanish outside a finite range.  Hence $C$ has finite
$p$-primary integral homology in finitely many degrees.  On
reduced mod-$p$ homology the map $p_Y$ is zero, so the mod-$p$
cofiber sequence injects $\widetilde H_k(Y;\mathbb F_p)$ into
$\widetilde H_k(C;\mathbb F_p)$.  Thus $C$ is noncontractible.

By \citet[Proposition~4C.1 and its proof, pp.~429--430]{A2Hatcher},
a simply connected CW complex with finitely generated homology
has a CW model with one generator cell in degree $k$ and one
relator cell in degree $k+1$ for each cyclic torsion summand of
$H_k$, and one cell for each free generator.  Here the homology
is finite and supported in finitely many degrees, so this model
$K\simeq C$ has finitely many cells.  It is the required finite
primary complex.  This argument allows divisible primary groups
in the original homology of $X$; no integral finiteness assumption
on $X$ was imposed.

\par\medskip\noindent\textbf{The finite-addition theorem and the exact periodicity input.}\quad
Write $W=\min\operatorname{Rev}(BC_p\vee V\to*)$.
Both generators are $\mathbb Z[1/p]$-acyclic, so $W\subseteq U_p$.
Replacing $V$ by a sufficiently high suspension does not change
$W$.  We may therefore assume that $V$ is a simply connected
suspension.  Its nonzero finite suspension spectrum has some
finite $p$-type $n+1\geq1$; here the finite-type classification
is the already rechecked input of \citet{AHopkinsSmith}.
Let $d$ be the lowest degree with $\pi_dV\ne0$, and choose an
odd $m\geq3$ with $m-1\geq d$.  The primary homology condition
also makes $V$ equivalent to its $p$-localization: that localization
is an integral homology equivalence between simply connected
spaces.  In what follows all Goodwillie stages are $p$-local.

Distinguish the ordinary nullification $P_V$ from the stages
$G_s=P_s^{\mathrm{Id}}(S^m_{(p)})$ of the Goodwillie tower.
\citet[Theorem~4.1, pp.~777--780]{EH5AroneMahowald1999}
prove that
\[
                  S^m_{(p)}\longrightarrow G_{p^n}
                  \tag{A4.7.4}
\]
is a $v_j$-periodic equivalence for every $0\leq j\leq n$,
where $v_0$ denotes rational homotopy.  This is the theorem
for odd spheres; no periodic convergence theorem for $V$ or
for arbitrary finite complexes is being used.

The comparison with $P_V$ is
\citet[Corollary~4.8, with the conventions of Section~4.3]{A3Bousfield2001Telescopic},
in the arbitrary finite suspension form stated explicitly by
\citet[Theorem~3.1, p.~8]{A4BarthelHeutsMeier2021}:
if a finite suspension $V$ has type $n+1$ and first nonzero
homotopy in degree $d$, a map between $d$-connected $p$-local
spaces is a $P_V$-equivalence precisely when it is a
$v_j$-equivalence for $0\leq j\leq n$.
The original Bousfield notation chooses a particular nullifier
of minimal connectivity.  The cited published 2021 formulation
is what permits the arbitrary $V$ used here.

We verify the connectivity hypothesis.  For $s>1$, the spectrum
of the $s$th layer is
\[
 E_s=\bigl(\operatorname{Map}_*(\Sigma K_s,
              \Sigma^\infty(S^m_{(p)})^{\wedge s})\bigr)_{h\Sigma_s},
 \qquad D_s(S^m_{(p)})\simeq\Omega^\infty E_s,
 \tag{A4.7.5}
\]
by \citet[Theorem~2.1]{EH5AroneMahowald1999}.
The partition complex $\Sigma K_s$ has dimension $s-1$.
Finite cotensoring therefore puts $E_s$ in spectrum degrees
$\geq sm-s+1\geq m$; homotopy orbits preserve this lower bound.
Starting with $G_1=QS^m_{(p)}$, the layer fiber sequences show
that each $G_s$ is $(m-1)$-connected and $p$-local.  Thus
\textup{(A4.7.4)} is a $P_V$-equivalence and belongs to $W$.
Ordinary nullification equivalences are contained in $W$ because
the cone-attachment units are generated by $V\to*$ and its
nonnegative suspensions, all inverted by $W$.

\par\medskip\noindent\textbf{Removing the finite tower and producing a Moore space.}\quad
For $s>1$, \citet[Proposition~3.1, pp.~758--759]{EH5AroneMahowald1999}
show that $E_s$ is rationally zero.  Since $E_s$ is $p$-local,
$E_s[1/p]=0$.  It is connected by the preceding lower bound.
The already rechecked \citet[Theorem~3.2]{ALevy} gives
\[
                         P_{BC_p}\Omega^\infty E_s\simeq *.
                         \tag{A4.7.6}
\]
A fibration with connected $A$-acyclic fibers is a
$P_A$-equivalence: for every $A$-null target $T$, the natural
constant maps $T\to\Map(F_b,T)$ are equivalences for the
fibers $F_b$, and taking their homotopy limit over the base
identifies maps from the total space with maps from the base.
Consequently the finitely many projections
$G_s\to G_{s-1}$ needed here are $P_{BC_p}$-equivalences.

The rational Hurewicz map
$QS^m_{(p)}\to K(\mathbb Q,m)$ is likewise a
$P_{BC_p}$-equivalence.  Its fiber is $\Omega^\infty F$, where
\[
 F=\operatorname{fib}(\Sigma^m\mathbb S_{(p)}
                              \longrightarrow\Sigma^mH\mathbb Q).
 \tag{A4.7.7}
\]
Here $F[1/p]=0$ and $F$ is concentrated in degrees
$\geq m-1\geq2$, so the same theorem of Levy applies.
Composing these equivalences gives a natural $W$-equivalence
$r:S^m_{(p)}\to K(\mathbb Q,m)$.  Naturality for a degree-$p$
map yields $r\circ p\simeq p\circ r$, where multiplication
by $p$ on $K(\mathbb Q,m)$ is an equivalence.  Two out of
three puts the degree-$p$ map of $S^m_{(p)}$ in $W$.

Its cofiber is homotopy equivalent to the finite Moore space
$M(\mathbb Z/p,m)$: the comparison from the integral
degree-$p$ cofiber is an integral homology equivalence of
simply connected spaces.  The reversible cofiber criterion,
followed by desuspension, makes $M=M(\mathbb Z/p,2)$
$W$-aspherical.

For completeness, an elementary nullification argument now
finishes without an additional Serre-class theorem.  Let $X$
be simply connected and primary, and put $Z=P_MX$.
The cone-attachment construction keeps $Z$ simply connected
and primary.  Applying pointed mapping spaces to
$S^2\xrightarrow{p}S^2\to M$ shows that
$\Map_*(M,Z)\simeq *$ is the homotopy fiber of
$p:\Omega^2 Z\to\Omega^2 Z$.  In particular multiplication
by $p$ is injective on $\pi_2Z$ and is an isomorphism on
$\pi_rZ$ for $r\geq3$.  If $Z$ were noncontractible, let
$r\geq2$ be its first nonzero homotopy degree.  Hurewicz
identifies $\pi_rZ$ with the nonzero primary group $H_r(Z;\mathbb Z)$.
Every nonzero primary torsion group has a nonzero element
annihilated by $p$, a contradiction.  Hence $P_MX\simeq *$,
and $X$ is $W$-aspherical.  The suspension of any connected
primary space has this form, so reversibility covers all
connected primary spaces.  These are exactly the $U_p$-aspherical
spaces.  The cofiber criterion gives $U_p\subseteq W$, proving
\textup{(A4.7.1)}.

\par\medskip\noindent\textbf{The H-space endpoint dichotomy.}\quad
Put $Z=P_{BC_p}H$.  The preservation argument of A4.5 makes
$Z$ simply connected and primary.  By the published
\citet[Theorem~7.2]{A4CCSDeconstruct2007}, its total mod-$p$
cohomology, and hence homology, is finite dimensional.
If $Z\simeq *$, the original H-space is ordinarily
$BC_p$-acyclic, giving the first line of \textup{(A4.7.2)}.
Otherwise its reduced mod-$p$ homology is nonzero: its vanishing
would imply $H_k(Z;\mathbb Z)[p]=0$ in every degree by universal
coefficients, hence integral acyclicity by primary torsion,
and then contractibility by simple connectivity.

The unit $H\to Z$ belongs to
$W_H=\min\operatorname{Rev}(BC_p\vee H\to*)$, so $Z$ is
$W_H$-aspherical.  The first assertion supplies a nonzero finite
primary $W_H$-aspherical complex.  Assertion~2 now gives
$W_H=U_p$.  Finally, for any set or class of the stated
H-space generators, either all are ordinarily $BC_p$-acyclic
or one is not.  The two alternatives just proved give the
two endpoints respectively.
\end{proof}

\paragraph{Scope of this presentation restriction.}
A4.7 does not assert that arbitrary localizers in $[W_0,U_p]$
admit presentations by H-spaces with finitely generated
Steenrod-algebra cohomology.  Degreewise finite type is not
substituted for that hypothesis.  Nor does passage to arbitrary
infinite homotopy colimits preserve the asserted finite total
cohomology.  The unrestricted third-element and full-fiber
set-size questions therefore remain open.  The finite-addition
assertion was supplied background; its fresh primary-source
audit is not claimed as a new solution of either question.

\begin{claim}{A4.8}{Proved}{Finite iterated-loop replacement preserves reversible generation}
\label{A48}
Let $W$ be a reversible fundamental localizer, let $X$ be a
connected well-pointed CW homotopy type, and let $1\leq m<\infty$.
Then
\[
 X\in\mathcal A_W
 \quad\Longleftrightarrow\quad
 \Omega^m\Sigma^mX\in\mathcal A_W,
 \qquad
 \min\operatorname{Rev}(X\to*)
   =\min\operatorname{Rev}(\Omega^m\Sigma^mX\to*).
 \tag{A4.8.1}
\]
Consequently every reversible class in the primary interval
$[W_0,U_p]$ admits a presentation over $W_0$ by simply connected
primary double-loop spaces: a connected generator $X$ may be
replaced by $\Omega^2\Sigma^3X$.  This applies equally to set
presentations and to class presentations in the fixed outer universe.
\end{claim}
\begin{proof}
Use the actual unstable filtration of the free little-$m$-cubes
space $C_mX$.  By \citet[Proposition~2.6(i)--(ii)]{A3May1972},
its successive pairs are NDR-pairs, $F_0C_mX=*$, and
\[
 F_rC_mX/F_{r-1}C_mX
       \cong \mathcal C_m(r)_+\wedge_{\Sigma_r}X^{\wedge r}.
 \tag{A4.8.2}
\]
The proof of that proposition constructs the NDR homotopies from
the $\Sigma_r$-equivariant NDR pair for the fat wedge in $X^r$;
see also its Appendix~A.4.  Thus these are homotopy cofiber
sequences, rather than only stable splitting terms.

Suppose $X\in\mathcal A_W$.  Smash closure puts
$\mathcal C_m(r)_+\wedge X^{\wedge r}$ in $\mathcal A_W$.
The symmetric group acts freely on the little-cubes factor
\citep[Corollary~4.6]{A3May1972}, so its based orbit in
\textup{(A4.8.2)} agrees with its derived based orbit.
Explicitly, compare the pairs
\[
 E\Sigma_r\times_{\Sigma_r}
       (\mathcal C_m(r)\times X^{\wedge r},\,
        \mathcal C_m(r)\times\{*\})
 \longrightarrow
 (\mathcal C_m(r)\times X^{\wedge r},\,
        \mathcal C_m(r)\times\{*\})/\Sigma_r.
\]
Both maps of spaces in these pairs are equivalences, by the free
configuration action, and the based orbits are their homotopy
cofibers.  The derived based orbit is a pointed homotopy colimit
of objects in $\mathcal A_W$ and hence lies there.  Cofiber
extension closure inductively gives $F_rC_mX\in\mathcal A_W$.
The filtration inclusions are cofibrations, so its union is the
sequential homotopy colimit and is also in $\mathcal A_W$.
The connected approximation theorem
\citep[Theorem~6.1]{A3May1972} gives
$C_mX\simeq\Omega^m\Sigma^mX$.

Conversely the suspension--loop adjunction triangle exhibits
$\Sigma^mX$ as a homotopy retract of
$\Sigma^m\Omega^m\Sigma^mX$.  Retract closure and $m$ applications
of reversible desuspension recover $X\in\mathcal A_W$.
The assertion for least reversible closures follows by applying
this equivalence to every $W$.

For the presentation statement first suspend a connected primary
$X$ once and use $m=2$.  The space $\Omega^2\Sigma^3X$ is simply
connected, and \textup{(A4.8.1)} applied also to $U_p$ shows that
it remains primary.  Replacing each generator in this way preserves
the whole reversible closure, without taking an impermissible wedge
of a proper class of objects.
\end{proof}

\paragraph{What the replacement does not supply.}
The loop-on-H presentation obstruction is removed.  Finite generation
of the replacement's cohomology as an unstable algebra is not proved
and does not hold in general.  A finite primary Moore generator, for
example, yields a double-loop generator for $U_p$ under A4.8;
whereas the CCS--Hemmi argument in A34.6 makes every simply connected
primary loop-on-H space with finitely generated unstable cohomology
ordinarily $BC_p$-acyclic.  No uniform cardinal bound on arbitrary
presentations follows.  The converse in \textup{(A4.8.1)} uses
finite $m$ and is not extended to $QX=\Omega^\infty\Sigma^\infty X$.
